\chapter*{Bảng viết tắt}
\addcontentsline{toc}{chapter}{Bảng viết tắt}

\begin{longtable}{|C{0.10\linewidth}|C{0.20\linewidth}|L{0.60\linewidth}|}
\caption*{Danh mục các từ viết tắt sử dụng trong báo cáo}\\
\hline
\textbf{STT} & \textbf{Từ viết tắt} & \textbf{Ý nghĩa} \\\hline
\endfirsthead
\hline
\textbf{STT} & \textbf{Từ viết tắt} & \textbf{Ý nghĩa} \\\hline
\endhead
1 & AI & Artificial Intelligence -- trí tuệ nhân tạo, lĩnh vực nghiên cứu và phát triển các hệ thống có khả năng thực hiện những nhiệm vụ thường cần đến trí thông minh của con người. \\\hline
2 & LLM & Large Language Model -- mô hình ngôn ngữ lớn, mô hình học sâu được huấn luyện trên lượng dữ liệu văn bản rất lớn để hiểu và sinh ngôn ngữ tự nhiên. \\\hline
3 & RAG & Retrieval-Augmented Generation -- truy xuất tăng cường tạo sinh, kỹ thuật kết hợp truy xuất tài liệu liên quan với mô hình sinh văn bản để tăng độ chính xác của phản hồi. \\\hline
4 & NLP & Natural Language Processing -- xử lý ngôn ngữ tự nhiên, lĩnh vực nghiên cứu cách máy tính phân tích, hiểu và tạo ra ngôn ngữ của con người. \\\hline
5 & OCR & Optical Character Recognition -- nhận diện ký tự quang học, kỹ thuật trích xuất văn bản từ hình ảnh (ví dụ tài liệu PDF dạng ảnh quét). \\\hline
6 & TTS & Text-to-Speech -- chuyển văn bản thành giọng nói, kỹ thuật tổng hợp tiếng nói từ dữ liệu văn bản để hỗ trợ nghe nội dung. \\\hline
7 & UI & User Interface -- giao diện người dùng, lớp thành phần mà người dùng trực tiếp thao tác để làm việc với hệ thống. \\\hline
8 & UX & User Experience -- trải nghiệm người dùng, tổng thể cảm nhận và mức độ thuận tiện khi người dùng tương tác với hệ thống. \\\hline
9 & API & Application Programming Interface -- giao diện lập trình ứng dụng, tập hợp quy tắc cho phép các thành phần phần mềm trao đổi dữ liệu với nhau. \\\hline
10 & PDF & Portable Document Format -- định dạng tài liệu điện tử giữ ổn định bố cục khi chia sẻ hoặc hiển thị trên nhiều thiết bị. \\\hline
11 & DOCX & Office Open XML Document -- định dạng tài liệu văn bản của Microsoft Word, thường dùng để soạn thảo và trao đổi nội dung có cấu trúc. \\\hline
12 & JWT & JSON Web Token -- định dạng token nhỏ gọn, tự chứa thông tin, dùng để xác thực và duy trì phiên đăng nhập không trạng thái (stateless). \\\hline
13 & SSE & Server-Sent Events -- giao thức truyền dữ liệu một chiều từ máy chủ đến trình duyệt qua một kết nối HTTP duy trì liên tục, dùng để truyền phản hồi AI theo luồng. \\\hline
14 & IDOR & Insecure Direct Object Reference -- lỗ hổng bảo mật xảy ra khi hệ thống cho phép truy cập trực tiếp vào tài nguyên chỉ dựa trên định danh, không kiểm tra quyền sở hữu. \\\hline
15 & HNSW & Hierarchical Navigable Small World -- thuật toán đồ thị dùng để tìm kiếm xấp xỉ láng giềng gần nhất (ANN) hiệu quả trong không gian vector nhiều chiều. \\\hline
16 & PostgreSQL & PostgreSQL -- hệ quản trị cơ sở dữ liệu quan hệ mã nguồn mở, dùng để lưu trữ dữ liệu hệ thống và lịch sử hội thoại. \\\hline
17 & pgvector & pgvector -- phần mở rộng của PostgreSQL cho phép lưu trữ, lập chỉ mục và truy vấn vector nhúng ngay trong cơ sở dữ liệu. \\\hline
18 & UUID & Universally Unique Identifier -- mã định danh duy nhất toàn cục, dùng để phân biệt các thực thể như tài liệu, phiên và tin nhắn. \\\hline
\end{longtable}
